% ============================================================
%  PART IV  --  My Data
% ============================================================
\part{My Data}

% ── Chapter 12: My Data Overview ──────────────────────────────
\chapter{Mis datos... Sinopsis}
\label{ch:mydata}

Mis datos son el administrador de recursos integrado de FairWindSK.
Proporciona una interfaz dedicada y optimizada para todos los datos de navegación almacenados en el Signal~K
servidor: waypoints, rutas, regiones, notas, gráficos, historial y archivos.

Open My Data by tapping the \key{My Data} icon (≡) in the Bottom Bar, or by tapping its tile on the
lanzador.

\screenshot{mydata-main-tabs}{Mis datos --- tab bar mostrando las siete pestañas}{}

\begin{notebox}
Todos los recursos en Mis datos se almacenan en el servidor Signal~K, no en el dispositivo FairWindSK.
Son accesibles desde cualquier dispositivo conectado al mismo servidor Signal~K simultáneamente.
\end{notebox}

\section{My Data Tab Overview}

Mis datos se organizan en siete pestañas.
Cada pestaña carga lazily --- se crea sólo la primera vez que se toca, manteniendo el arranque rápido.

\rowcolors{2}{LtGray}{white}
\begin{longtable}{C{0.7cm} L{2.8cm} L{10.2cm}}
\toprule
\rowcolor{Navy}
{\sffamily\bfseries\color{white} Tab} &
{\sffamily\bfseries\color{white} Name} &
{\sffamily\bfseries\color{white} Contents} \\
\midrule
0 & Waypoints & Named positions: anchorages, hazards, marinas, fishing marks, race marks. \\
1 & Routes & Ordered sequences of waypoints forming planned passages. \\
2 & Regions & Geographic boundary areas --- SAR search regions, exclusion zones, anchorage limits. \\
3 & Notes & Free-text annotations attached to positions. \\
4 & Charts & Chart tile source catalogue for installed chart providers. \\
5 & Tracks & Historical GPS track --- the vessel's recorded position history. \\
6 & Files & Direct filesystem browser for files accessible through Signal~K. \\
\bottomrule
\end{longtable}

\section{La barra de herramientas común}

Tabs 0--4 (Waypoints, Routes, Regiones, Notes, Charts) comparten un diseño de interfaz común:

\begin{itemize}
  \item A \textbf{title label} showing the resource type.
  \item A \textbf{search field} for real-time filtering of the list. Searching is case-insensitive and checks all visible columns simultaneously.
  \item A \textbf{scrollable table} showing all resources from the Signal~K server. Rows are 64\,px tall for comfortable touch selection. Alternating row colours aid readability. Column headers are clickable for sorting.
  \item A \textbf{toolbar} of icon buttons:
\end{itemize}

\rowcolors{2}{LtGray}{white}
\begin{longtable}{L{3cm} L{10.7cm}}
\toprule
\rowcolor{Navy}
{\sffamily\bfseries\color{white} Button} &
{\sffamily\bfseries\color{white} Action} \\
\midrule
Open ($\rightarrow$) & Opens the selected resource --- shows its details, opens an embedded map preview (Waypoints), or views the attachment (Notes). Accent-coloured; primary action. \\
Edit & Opens the resource editor drawer for the selected resource. \\
New (+) & Opens the resource editor drawer to create a new resource. \\
Delete & Deletes the selected resource after confirmation. Cannot be undone. \\
Import & Imports resources from a GeoJSON or JSON file. File picker opens in the Bottom Bar drawer. \\
Export & Exports the selected resource (or all resources if none selected) to a GeoJSON file. \\
Refresh & Forces a reload of all resources from the Signal~K server. \\
\bottomrule
\end{longtable}

\begin{notebox}
En Raspberry Pi OS (ARM Linux), la barra de herramientas utiliza botones solo de texto para evitar la reproducción de iconos
cuestiones en esa plataforma.
El diseño de botones y la funcionalidad son idénticos; sólo la presentación visual difiere.
\end{notebox}

% ── Chapter 13: Waypoints ─────────────────────────────────────
\chapter{Waypoints}
\label{ch:waypoints}

Waypoints son nombrados posiciones almacenadas en el servidor Signal~K.
Se utilizan para anclajes, peligros, puertos deportivos, paradas de combustible, marcas de pesca, marcas de carreras, y cualquier
otra posición a la que quieres grabar y navegar.

\screenshot{mydata-waypoints-list}{Mis datos... Ficha Waypoints, vista lista con nombre, descripción, tipo, lat/lon, timestamp}{}

\section{Waypoints Table Columns}

La tabla waypoints muestra seis columnas:

\rowcolors{2}{LtGray}{white}
\begin{longtable}{L{2.4cm} L{11.3cm}}
\toprule
\rowcolor{Navy}
{\sffamily\bfseries\color{white} Column} &
{\sffamily\bfseries\color{white} Description} \\
\midrule
Name & Display name of the waypoint, taken from the \code{name} field or from \code{feature.properties.name} in the GeoJSON resource. \\
Description & Free-text description of the waypoint. \\
Type & The waypoint type tag, e.g.\ \code{alarm-mob} (POB waypoints created by FairWindSK), \code{anchor}, \code{marina}, or any custom type. POB waypoints are visually distinguishable by their type. \\
Latitude & Latitude in your configured coordinate format (right-aligned). \\
Longitude & Longitude in your configured coordinate format (right-aligned). \\
Timestamp & ISO 8601 creation or last-update timestamp. \\
\bottomrule
\end{longtable}

\section{Searching Waypoints}

El campo de búsqueda por encima de la tabla filtra waypoints por cualquier columna visible.
Escriba parte de un nombre, descripción, tipo o valor de coordenadas.
La tabla actualiza en tiempo real como escribes.

\section{Crear un waypoint}

Tap \key{New} to open the Waypoint Editor drawer.


\subsection{Editor Fields}

\rowcolors{2}{LtGray}{white}
\begin{longtable}{L{3cm} L{10.7cm}}
\toprule
\rowcolor{Navy}
{\sffamily\bfseries\color{white} Field} &
{\sffamily\bfseries\color{white} Description} \\
\midrule
Name & The display name shown in the table, Top Bar waypoint widget, and chart plotter. \\
Description & Free-text notes about this position. \\
Type & Optional type tag. Leave blank for a generic waypoint. Use \code{anchor} for anchoring positions, \code{hazard} for obstructions, etc. \\
Latitude & Degrees with decimal minutes or decimal degrees, depending on the active coordinate format. Range: $-90$ to $+90$. \\
Longitude & Range: $-180$ to $+180$. \\
Altitude & Altitude in metres (optional). Leave as 0 if not relevant. \\
\bottomrule
\end{longtable}

After filling in the fields, tap \key{OK} in the editor.
El waypoint se crea inmediatamente en el servidor Signal~K y aparece en la tabla.

\section{Edición de un waypoint}

Select a waypoint in the table, then tap \key{Edit}.
El mismo cajón del editor se abre con los valores existentes pre-filled.
Modify any field and tap \key{OK}.
El cambio está escrito inmediatamente a Signal~K.

\section{Eliminar un waypoint}

Select a waypoint and tap \key{Delete}.
Aparece un cajón de confirmación.
Tap \key{Delete} to confirm.
El waypoint se elimina del servidor Signal~K permanentemente.

\begin{warningbox}
La eliminación no se puede deshacer.
POB waypoints (type \code{alarm-mob}) should generally be retained for accident reporting and debrief,
no eliminado.
\end{warningbox}

\section{Abrir un waypoint --- la vista de detalle}

Select a waypoint and tap \key{Open} (or double-tap the row).
La vista Waypoint Detail abre.


La vista detallada muestra:

\begin{itemize}
  \item Full name, description, type, and formatted coordinates.
  \item An embedded Freeboard-SK map preview centred on the waypoint's position, if Freeboard-SK is installed on the connected Signal~K server. The map uses a JavaScript function to compute the view bounding box from the waypoint coordinates and zoom level.
  \item A \key{Navigate} button to set the waypoint as the active navigation destination.
\end{itemize}

\begin{notebox}
La vista previa integrada Freeboard-SK está deshabilitada en Raspberry Pi OS (ARM Linux) se construye para evitar
problemas de rendimiento con casos WebEngine simultáneos en esa plataforma.
Las coordenadas y descripciones siempre se muestran en forma de texto independientemente.
\end{notebox}

\section{Navegando a un Waypoint}

From the detail view, tap \key{Navigate}.
FairWindSK calls \path{navigateToWaypoint} on the Signal~K client, which writes the waypoint's
resource path (e.g.\ \path{/resources/waypoints/\{uuid\}}) as the active navigation destination on
Signal~K.
Inmediatamente:

\begin{itemize}
  \item The Top Bar BTW, DTW, TTG, and ETA widgets begin showing values relative to this waypoint.
  \item Freeboard-SK and any other navigation application subscribed to \skpath{navigation.course} show the destination.
  \item The POB bearing and distance displays (when the POB bar is active) use this same destination.
\end{itemize}

\section{Importing Waypoints}

Tap \key{Import}. A file picker opens in the Bottom Bar drawer.
Select a GeoJSON file containing \code{Feature} objects with \code{Point} geometry.

FairWindSK:
\begin{enumerate}
  \item Reads the GeoJSON file.
  \item Extracts each \code{Point} feature as a waypoint.
  \item Creates or updates each waypoint on the Signal~K server (update if a matching UUID is found; create if not).
  \item Reports the count of imported waypoints in an informational drawer.
  \item Selects the first imported waypoint in the table.
\end{enumerate}

The import file filter accepts \code{.geojson} and \code{.json} files.

\section{Exporting Waypoints}

Tap \key{Export}. If a waypoint is selected, only that waypoint is exported.
Si no se selecciona nada, se exportan todos los puntos.
Se abre un selector de archivo. Elija un nombre de archivo y ubicación.
The output is a valid GeoJSON \code{FeatureCollection} with \code{Point} geometry features.

% ── Chapter 14: Routes ────────────────────────────────────────
\chapter{Rutas}
\label{ch:routes}

Las rutas se ordenan secuencias de waypoints formando pasajes planeados.
Normalmente se crean en un plotter de gráficos (Freeboard-SK) y se almacenan en el servidor Signal~K.

\screenshot{mydata-routes-list}{Mis datos... Ficha de ruta, vista de lista con nombre, tipo geometría, conteo de puntos, timetamp}{}

\section{Rutas Columnas de mesa}

\rowcolors{2}{LtGray}{white}
\begin{longtable}{L{2.4cm} L{11.3cm}}
\toprule
\rowcolor{Navy}
{\sffamily\bfseries\color{white} Column} &
{\sffamily\bfseries\color{white} Description} \\
\midrule
Name & Display name of the route. \\
Description & Free-text description or summary. \\
Geometry & The GeoJSON geometry type --- typically \code{LineString} for simple routes, \code{MultiLineString} for multi-leg routes with gaps, or \code{sar-spiral} / \code{sar-parallel} for SAR patterns created by the POB feature. \\
Points & The number of coordinate points in the route geometry. Displayed centred. \\
Timestamp & ISO 8601 creation or last-update timestamp. \\
\bottomrule
\end{longtable}

\section{Abrir una ruta --- la vista Detalle}

Select a route and tap \key{Open}.


La vista detallada muestra:
\begin{itemize}
  \item Name, description, geometry type, and point count.
  \item A GeoJSON geometry preview if available.
  \item An \key{Activate} button to set the route's first waypoint as the active navigation destination.
\end{itemize}

\section{Creación y edición de rutas}

\subsection{Crear una ruta}

Tap \key{New}. The Route Editor drawer opens.

\rowcolors{2}{LtGray}{white}
\begin{longtable}{L{3cm} L{10.7cm}}
\toprule
\rowcolor{Navy}
{\sffamily\bfseries\color{white} Field} &
{\sffamily\bfseries\color{white} Description} \\
\midrule
Name & Display name of the route. \\
Description & Free-text notes. \\
Coordinates (JSON) & The GeoJSON coordinate array for the route geometry. Enter as a JSON array of \code{[longitude, latitude]} pairs. \\
Geometry (JSON) & Advanced: edit the full GeoJSON geometry object if needed. \\
Properties (JSON) & Advanced: additional GeoJSON feature properties. \\
\bottomrule
\end{longtable}

\subsection{Activando una ruta}

From the route detail view, tap \key{Activate}.
FairWindSK establece el primer punto de la ruta como el destino activo de navegación Signal~K.
The Autopilot Route mode (Chapter~\ref{ch:autopilot}) will then follow the active route automatically.

\section{Rutas de importación y exportación}

Import and export use the same GeoJSON \code{FeatureCollection} format as Waypoints.
El tipo de geometría en el GeoJSON determina cómo se almacena la ruta en Signal~K.

\begin{tipbox}
SAR search routes created by the POB feature (Chapter~\ref{ch:pob}) appear in the Routes tab
automáticamente después de la creación.
They have geometry types \code{sar-spiral} and \code{sar-parallel} and can be exported as GeoJSON
para los registros de viaje.
\end{tipbox}

% ── Chapter 15: Regions ───────────────────────────────────────
\chapter{Regiones}
\label{ch:regions}

Las regiones son áreas fronterizas geográficas definidas como poligones GeoJSON.
Se utilizan para áreas de búsqueda SAR (automáticamente creadas por la función POB), área protegida marina
límites, zonas de exclusión, límites de anclaje y zonas de pesca.

\screenshot{mydata-regions-list}{Mis datos... Ficha de regiones, vista de lista con nombre, tipo de geometría, timetamp}{}

\section{Regiones Columnas Tabla}

\rowcolors{2}{LtGray}{white}
\begin{longtable}{L{2.4cm} L{11.3cm}}
\toprule
\rowcolor{Navy}
{\sffamily\bfseries\color{white} Column} &
{\sffamily\bfseries\color{white} Description} \\
\midrule
Name & Display name of the region. \\
Description & Free-text description. \\
Geometry & The GeoJSON geometry type --- typically \code{Polygon} or \code{MultiPolygon}. SAR regions use \code{Polygon}. \\
Timestamp & ISO 8601 creation or last-update timestamp. \\
\bottomrule
\end{longtable}

\section{Creación de una región}

Tap \key{New}. The Region Editor drawer opens.

\rowcolors{2}{LtGray}{white}
\begin{longtable}{L{3cm} L{10.7cm}}
\toprule
\rowcolor{Navy}
{\sffamily\bfseries\color{white} Field} &
{\sffamily\bfseries\color{white} Description} \\
\midrule
Name & Display name. \\
Description & Free-text notes. \\
Coordinates (JSON) & The GeoJSON coordinate array for the polygon. Polygons require an outer ring (first and last coordinate identical to close the ring). \\
Geometry (JSON) & Advanced: full GeoJSON geometry object. \\
Properties (JSON) & Advanced: additional GeoJSON feature properties. \\
\bottomrule
\end{longtable}

\begin{notebox}
Áreas de búsqueda SAR creadas por la función POB aparecen en la pestaña Regiones automáticamente.
Their \code{properties} include the \code{pobUuid} linking them to the originating POB incident.
\end{notebox}

% ── Chapter 16: Notes ─────────────────────────────────────────
\chapter{Notas}
\label{ch:notes}

Las notas son anotaciones de texto libre adjuntas a posiciones geográficas.
Funcionan como notas pegajosas en el gráfico --- útil para grabar consejos de pilotaje, peligros locales no
en el gráfico, detalles del puerto o observaciones de la tripulación.

\screenshot{mydata-notes-list}{Mis datos... Notas, vista de lista con título, descripción, href, tipo MIME, timetamp}{}

\section{Notas Columnas de tabla}

Las notas tienen un conjunto de columnas que reflejan su naturaleza orientada hacia el apego:

\rowcolors{2}{LtGray}{white}
\begin{longtable}{L{2.4cm} L{11.3cm}}
\toprule
\rowcolor{Navy}
{\sffamily\bfseries\color{white} Column} &
{\sffamily\bfseries\color{white} Description} \\
\midrule
Title & The display name of the note (the first column is labelled \textit{Title} for notes, not \textit{Name}, to reflect the note-taking context). \\
Description & Summary text. For imported text files, this is set to the first 120 characters of the file content. \\
Href & An optional file path or URL linking the note to an attachment --- a PDF pilot book page, an image, a text file. When present, tapping \key{Open} opens the attachment directly in the built-in viewer. \\
MIME Type & The MIME type of the attachment, e.g.\ \code{application/pdf}, \code{image/jpeg}, \code{text/plain}. Used to route the file to the correct viewer. \\
Timestamp & ISO 8601 creation or last-update timestamp. \\
\bottomrule
\end{longtable}

\section{Crear una nota}

Tap \key{New}. The Note Editor drawer opens.


\rowcolors{2}{LtGray}{white}
\begin{longtable}{L{3cm} L{10.7cm}}
\toprule
\rowcolor{Navy}
{\sffamily\bfseries\color{white} Field} &
{\sffamily\bfseries\color{white} Description} \\
\midrule
Title & The note title (displayed in the \textit{Title} column). \\
Description & Main text content of the note. \\
Href & Optional path or URL to an attachment file. Enter an absolute path (e.g.\ \path{/home/pi/charts/harbour.pdf}) or a URL. \\
MIME Type & Optional MIME type for the attachment. FairWindSK uses this to select the correct viewer. If left blank, FairWindSK detects the type from the file extension. \\
Has position & Check this box to attach a geographic position to the note. \\
Latitude & Latitude when \textit{Has position} is checked. \\
Longitude & Longitude when \textit{Has position} is checked. \\
\bottomrule
\end{longtable}

\section{Apertura de una nota}

\subsection{Notas sin un adjunto}

Tapping \key{Open} opens the note editor in read-only mode, showing the title, description, and position.

\subsection{Notas con un acoplamiento (conjunto de href)}

Tapping \key{Open} immediately opens the attachment:

\begin{itemize}
  \item \textbf{Local file that exists:} opens in the built-in FileViewer (PDF, image, or text viewer depending on the MIME type).
  \item \textbf{HTTP/HTTPS URL:} opens in an embedded web view inside the Bottom Bar horizontal drawer.
  \item \textbf{Local directory path:} shows a message directing you to use My Data~> Files to browse the folder (single-window model prevents opening a separate file manager).
\end{itemize}

\section{Notas de importación}

La pestaña Notas acepta un conjunto más amplio de tipos de archivos de importación que otras pestañas.
Además de GeoJSON/JSON:

\begin{itemize}
  \item \textbf{Text files} (\code{.txt}, \code{.md}, \code{.html}): FairWindSK creates a note with the filename as the title and the first 120 characters of the file content as the description, setting the \code{href} to the absolute file path and detecting the MIME type automatically.
  \item \textbf{Images} (\code{.jpg}, \code{.jpeg}, \code{.png}): Creates a note with the filename as title and sets the \code{href} and \code{mimeType} accordingly.
  \item \textbf{PDF files}: Creates a note pointing to the PDF file via \code{href}.
\end{itemize}

The import filter text reads: \textit{Notes and attachments (*.geojson *.json *.txt *.md *.markdown *.html *.htm *.pdf *.jpg *.jpeg *.png)}.

% ── Chapter 17: Charts ────────────────────────────────────────
\chapter{Gráficos}
\label{ch:charts}

La pestaña Charts lista las fuentes disponibles de fichas de mapa registradas en el servidor Signal~K.
Gestiona el catálogo fuente de la gráfica --- no la pantalla de la gráfica real, que sucede dentro
Freeboard-SK u otro plotter gráfico que se ejecuta en el Área de Aplicación.

\screenshot{mydata-charts-list}{Mis datos... Ficha de gráficos, vista de lista con nombre, formato, identificador, URL, timetamp}{}

\section{Columnas Tabla Columnas}

\rowcolors{2}{LtGray}{white}
\begin{longtable}{L{2.4cm} L{11.3cm}}
\toprule
\rowcolor{Navy}
{\sffamily\bfseries\color{white} Column} &
{\sffamily\bfseries\color{white} Description} \\
\midrule
Name & Display name of the chart source. \\
Description & Free-text description. \\
Format & The chart tile format. Supported formats include: \code{mbtiles} (MBTiles SQLite database), \code{kap} (BSB/KAP raster charts), \code{gemf} (GEMF tile archive), \code{tilemapresource} (TileMapResource XML), \code{package} (ZIP chart package). \\
Identifier & The chart identifier used by Signal~K and chart plotters to reference this source. \\
Chart URL & The local filesystem path or URL where the chart tiles are served. \\
Timestamp & ISO 8601 registration timestamp. \\
\bottomrule
\end{longtable}

\section{Añadiendo una fuente de carga}

Tap \key{New}. The Chart Editor drawer opens.

\rowcolors{2}{LtGray}{white}
\begin{longtable}{L{3cm} L{10.7cm}}
\toprule
\rowcolor{Navy}
{\sffamily\bfseries\color{white} Field} &
{\sffamily\bfseries\color{white} Description} \\
\midrule
Name & Display name for this chart source. \\
Description & Free-text description. \\
Identifier & Unique identifier string used by chart plotters to reference this source. \\
Chart format & The format of the chart tiles. FairWindSK detects the format from the file extension when importing. \\
Chart URL & Path or URL to the chart tiles. For MBTiles, this is the absolute path to the \code{.mbtiles} file. \\
Tilemap URL & For \code{tilemapresource} format only: the path to the \code{tilemapresource.xml} file. The Chart URL is then set to the directory containing it. \\
Chart region & Optional: a GeoJSON bounding box or geometry describing the geographic coverage of this chart. \\
Chart scale & Optional: the chart scale denominator (e.g.\ \code{50000} for 1:50,000). \\
Chart layers & Optional: comma-separated list of layer names for tile sources that support multiple layers. \\
Chart bounds & Optional: GeoJSON coordinates describing the chart bounds. \\
\bottomrule
\end{longtable}

\section{Importar una fuente de carga}

Tap \key{Import}. Select a chart file.
FairWindSK detecta el formato de la gráfica desde la extensión de archivo y construye un objeto de recurso de la gráfica
automáticamente.
For \code{tilemapresource.xml} files, FairWindSK sets both the tile URL (the directory) and the
URL de tilemap (la ruta de archivo XML).

La importación crea o actualiza el registro gráfico en el servidor Signal~K.

\section{Exportar el catálogo de carga}

Tap \key{Export}. The chart catalogue is exported as a JSON manifest document with structure:

\begin{lstlisting}
{
  "type": "FairWindSKChartPackage",
  "charts": [ { ... }, { ... } ]
}
\end{lstlisting}

Este manifiesto puede ser importado en otra instalación de FairWindSK para reproducir el catálogo de gráficos
configuración (que describe las fuentes de gráficos, no los datos de la ficha propia).

% ── Chapter 18: Tracks ────────────────────────────────────────
\chapter{Pistas... Historia}
\label{ch:tracks}

La ficha Tracks muestra la pista histórica del GPS del buque --- la secuencia registrada de la posición
se arregla con el tiempo.
Las pistas se generan automáticamente por Signal~K cuando el plugin de pistas está habilitado.

\screenshot{mydata-tracks-list}{Mis datos... Tabla de pistas, tabla de puntos de pista con tiempo, lat/lon, altitud}{}

\section{Columnas de mesa de seguimiento}

La tabla de pistas muestra fijaciones individuales de posición (puntos) en orden cronológico:

\rowcolors{2}{LtGray}{white}
\begin{longtable}{L{2.4cm} L{11.3cm}}
\toprule
\rowcolor{Navy}
{\sffamily\bfseries\color{white} Column} &
{\sffamily\bfseries\color{white} Description} \\
\midrule
Time & The timestamp of the position fix, formatted in local time using the system locale. \\
Latitude & Latitude to 6 decimal places. \\
Longitude & Longitude to 6 decimal places. \\
Altitude & Altitude in metres (if available from the GPS fix; otherwise shown as 0). \\
\bottomrule
\end{longtable}

\section{Grabación de pistas en vivo}

Si Signal~K está proporcionando actualizaciones de posición activas, nuevos puntos de pista se adjuntan a la tabla en
tiempo real.
Los pergaminos de mesa para mostrar el punto más reciente automáticamente.

\section{Ver un punto de seguimiento}

Select a row and tap \key{Open}. A read-only dialog opens showing the timestamp, latitude, longitude,
y altitud para ese punto en detalle.

\subsection{Edición de un punto de seguimiento}

Select a row and tap \key{Edit}. The Track Point Editor opens.

\rowcolors{2}{LtGray}{white}
\begin{longtable}{L{3cm} L{10.7cm}}
\toprule
\rowcolor{Navy}
{\sffamily\bfseries\color{white} Field} &
{\sffamily\bfseries\color{white} Description} \\
\midrule
Timestamp & Date and time of the fix. Uses a date/time picker with up/down arrows. \\
Latitude & Latitude (8 decimal places, range $-90$ to $+90$). \\
Longitude & Longitude (8 decimal places, range $-180$ to $+180$). \\
Altitude & Altitude in metres. \\
\bottomrule
\end{longtable}

\section{Agregar un punto de seguimiento manualmente}

Tap \key{New}. The same editor opens with default values (current time and position 0,0).
Enter the correct values and tap \key{OK}.
El punto se inserta en la pista en orden cronológico.

\section{Eliminar un punto de pista}

Select a row and tap \key{Delete}. Confirm when prompted. The point is removed from the track.

\section{Exportar la pista}

Tap \key{Export} to save the complete track as a GeoJSON \code{LineString} or
\code{FeatureCollection}:

\begin{lstlisting}
{
  "type": "FeatureCollection",
  "features": [
    {
      "type": "Feature",
      "geometry": {
        "type": "LineString",
        "coordinates": [[lon, lat, alt], ...]
      }
    }
  ]
}
\end{lstlisting}

The default export filename is \path{tracks-history.geojson}.

\begin{tipbox}
Exportar la pista después de cada pasaje y almacenarla con el diario de viaje.
El archivo GeoJSON se puede importar en la mayoría de herramientas GIS, plotters de gráficos y software de planificación de viajes
para su examen y análisis.
\end{tipbox}

\section{Importar una pista}

Tap \key{Import} and select a GeoJSON file containing a \code{LineString} geometry.
FairWindSK extrae cada coordenadas en la línea como un punto de pista con un timetamp sintético,
los pasa a la pista actual en orden.

\section{Usos de la pestaña de pista}

\begin{itemize}
  \item \textbf{Passage review:} compare the actual track against the planned route. Identify deviations, traffic avoidance manoeuvres, and current effects.
  \item \textbf{Anchorage assessment:} see the swing circle after anchoring overnight.
  \item \textbf{Safety investigation:} reconstruct the vessel's path before and during an incident for formal reporting.
  \item \textbf{Logbook cross-reference:} correlate positions at specific times with logbook entries.
\end{itemize}

% ── Chapter 19: Files ─────────────────────────────────────────
\chapter{Archivos}
\label{ch:files}

La pestaña Archivos es un navegador directo para archivos y directorios accesibles en el dispositivo
Corriendo FairWindSK.
Proporciona operaciones de navegación, búsqueda, visualización y gestión de archivos (copia, corte, pasta, renombre,
borrar, nueva carpeta) sin dejar la interfaz FairWindSK.

\screenshot{mydata-files-browser}{Mis datos... Ficha de archivos, navegador del directorio con barra de herramientas}{}

\section{Interface Layout}

La pestaña Archivos se divide en tres áreas:

\begin{itemize}
  \item \textbf{Toolbar} at the top: navigation and action buttons, path bar, and search field.
  \item \textbf{File list} in the centre: the current directory contents as a tree view with columns for Name, Size, and Date. Row height is 58\,px for comfortable touch selection. Alternating row colours aid readability.
  \item \textbf{Actions row} below the toolbar: file operation buttons that appear when items are selected.
\end{itemize}

\section{Navegación}

\subsection{Directorio Navegación}

\rowcolors{2}{LtGray}{white}
\begin{longtable}{L{3cm} L{10.7cm}}
\toprule
\rowcolor{Navy}
{\sffamily\bfseries\color{white} Action} &
{\sffamily\bfseries\color{white} How to trigger} \\
\midrule
Enter a folder & Double-tap a folder row in the file list, or select and tap \key{Open}. \\
Go up one level & Tap \key{↑} (Up) button in the toolbar. \\
Navigate by path & Tap the path bar at the top of the toolbar, type an absolute path, and press Return. \\
Go home & Tap \key{Home} button. Returns to the home directory of the running user. \\
\bottomrule
\end{longtable}

\subsection{Clasificación de columnas}

Toque cualquier encabezado de columna para ordenar por esa columna.
La dirección del tipo se mueve entre ascender y descender en los golpes posteriores.
El tipo predeterminado es por Nombre, ascendente.

\section{Búsqueda de archivos}

\subsection{Búsqueda básica}

Tap \key{Search} or press Return in the search field.
FairWindSK busca recursivamente el árbol del directorio actual para archivos y subdirectorios cuyos nombres
contiene el texto de búsqueda.

Los resultados se muestran en una lista plana en la vista del archivo.
Toque un resultado para seleccionarlo; doble punción para abrir.
La búsqueda continúa hasta la raíz del árbol del directorio navegado.


\subsection{Opciones de búsqueda}

Tap \key{Options} (or the filter button) to expand the search options:

\rowcolors{2}{LtGray}{white}
\begin{longtable}{L{3.5cm} L{10.2cm}}
\toprule
\rowcolor{Navy}
{\sffamily\bfseries\color{white} Option} &
{\sffamily\bfseries\color{white} Description} \\
\midrule
Case Sensitive & When enabled, the search matches the exact capitalisation of the search text. Default: off (case-insensitive). \\
Search Hidden & When enabled, hidden files and directories (names beginning with \code{.}) are included in search results. Default: off. \\
Search System & When enabled, system files and directories (e.g.\ \code{/proc}, \code{/sys} on Linux) are included. Default: off. \\
\bottomrule
\end{longtable}

The search options button label updates to show a compact summary of active options (e.g.\
\textit{Options: case, hidden}).

\section{File Actions}

\subsection{Archivos de apertura}

Select a file and tap \key{Open}, or double-tap the row.

FairWindSK detecta el tipo de archivo mediante detección de tipo MIME y lo dirige al espectador adecuado:

\rowcolors{2}{LtGray}{white}
\begin{longtable}{L{4cm} L{9.7cm}}
\toprule
\rowcolor{Navy}
{\sffamily\bfseries\color{white} File type} &
{\sffamily\bfseries\color{white} Viewer} \\
\midrule
Images (\code{.jpg}, \code{.jpeg}, \code{.png}, \code{.gif}) & Built-in image viewer with pinch-to-zoom and pan. \\
PDF (\code{.pdf}) & Built-in PDF viewer, page by page. \\
Text files (\code{.txt}, \code{.md}, \code{.log}, \code{.html}) & Built-in text viewer, formatted as plain text or HTML. \\
Other file types & A message is shown; the file cannot be previewed but can be managed. \\
\bottomrule
\end{longtable}

Los directorios se abren en el navegador del directorio (navegando en ellos) en lugar de en un espectador.
Intento abrir una ruta de directorio desde un Notes href muestra un mensaje que le dirige a utilizar los archivos
pestaña para navegar en su lugar (principio modelo de ventanilla única).


\subsection{Recuperar archivos y carpetas}

Select an item and tap \key{Rename}.
Un diálogo de entrada se abre en el cajón horizontal Bottom Bar.
Enter the new name and tap \key{OK}.
El archivo o directorio se renombra inmediatamente en el sistema de archivos.

\subsection{Crear una nueva carpeta}

Tap \key{New Folder}.
Se abre un diálogo de entrada.
Enter the folder name and tap \key{OK}.
La nueva carpeta se crea en el directorio actual.

\subsection{Copiar, Cortar y Pegar}

\begin{enumerate}
  \item Select one or more files or folders.
  \item Tap \key{Copy} to mark them for copying, or \key{Cut} to mark them for moving.
  \item Navigate to the destination directory.
  \item Tap \key{Paste}.
\end{enumerate}

FairWindSK realiza controles de seguridad antes de las operaciones de pasta:

\begin{itemize}
  \item Copying a directory into itself or one of its subdirectories is blocked.
  \item Pasting into the source directory of a cut operation shows a warning.
  \item The self-paste safety check runs a self-test at startup to verify the detection logic is working correctly.
\end{itemize}

Las operaciones de copia y movimiento del directorio son recursivas --- todo el subárbol es copiado o movido.

\subsection{Eliminar archivos y carpetas}

Select one or more items and tap \key{Delete}.
Un diálogo de confirmación se abre en el cajón horizontal Bottom Bar.
Confirme su eliminación.

\begin{warningbox}
La eliminación de archivos es permanente.
No hay papel de reciclaje o deshacer para las operaciones de archivos.
Tenga especial cuidado al eliminar directorios --- todos los contenidos se eliminan recursivamente.
\end{warningbox}

\section{Columnas Lista de Archivo}

La lista de archivos muestra tres columnas:

\rowcolors{2}{LtGray}{white}
\begin{longtable}{L{2cm} L{11.7cm}}
\toprule
\rowcolor{Navy}
{\sffamily\bfseries\color{white} Column} &
{\sffamily\bfseries\color{white} Description} \\
\midrule
Name & Filename or directory name, with an icon indicating the file type (folder, image, document, etc.). \\
Size & File size in human-readable units (B, KB, MB, GB). Directories show the size of their immediate contents. \\
Date & Last modification timestamp. The \textit{Type} column (column index 2) is hidden by default. \\
\bottomrule
\end{longtable}

\section{Usos prácticos de la pestaña Archivos}

\begin{itemize}
  \item Browse pilot book excerpts stored as PDFs on the Raspberry Pi.
  \item View passage planning images (harbour photos, aerial views, approach charts).
  \item Access voyage log files or NMEA recordings.
  \item Manage chart tile archives before registering them in the Charts tab.
  \item Review diagnostic log files generated by FairWindSK (see Chapter~\ref{ch:settings-system}).
\end{itemize}
